% 土星（综述）
% license CCBYSA3
% type Wiki

本文根据 CC-BY-SA 协议转载翻译自维基百科\href{https://en.wikipedia.org/wiki/Saturn}{相关文章}。

\begin{figure}[ht]
\centering
\includegraphics[width=8cm]{./figures/13ddb48d3e7de073.png}
\caption{土星及其显著的行星环，由卡西尼轨道器拍摄\(^\text{[a]}\)} \label{fig_Saturn_1}
\end{figure}
土星是距离太阳第六颗行星，也是太阳系中仅次于木星的第二大行星。它是一颗气态巨行星，其平均半径约为地球的 9 倍。它的平均密度仅为地球的八分之一，但质量却超过地球的 95 倍。尽管土星几乎与木星一样大，但其质量还不到木星的三分之一。土星以 9.59 AU（14.34 亿公里）的距离绕太阳运行，轨道周期为 29.45 年。

土星内部被认为由一颗岩石核心组成，其外包裹着一层深厚的金属氢层，其上为液态氢和液态氦的中间层，最外层为气体外壳。土星呈现淡黄色调，是由于其上层大气中的氨晶体所致。金属氢层中的电流被认为会产生土星的行星磁场，该磁场强度弱于地球，但由于土星体积更大，其磁矩是地球的 580 倍。土星磁场的强度约为木星的二十分之一\(^\text{[27]}\)。土星外层大气整体较为平淡、缺乏对比度，但亦可能出现长期存在的结构。土星的风速可达每小时 1,800 公里（1,100 英里）。

这颗行星拥有明亮且庞大的光环系统，主要由冰粒构成，另含少量岩石碎屑与尘埃。至少有 274 颗卫星绕火星运行，其中 63 颗已有正式命名；这一数字不包含光环中的数百颗小卫星体。土卫六（Titan）是土星最大的卫星，也是太阳系中第二大的卫星，其体积大于（但质量小于）水星，并且是太阳系中唯一拥有稠密大气层的卫星\(^\text{[28]}\)。
\subsection{名称与符号}
土星以罗马财富与农业之神命名，该神是木星（Jupiter）之父。其天文符号（♄）可追溯至希腊的《奥克绪林库斯纸草文稿》（Oxyrhynchus Papyri），在其中该符号可见为希腊字母 κ–ρ（kappa–rho）的连写并带有一条水平短线，作为 Κρονος（Cronus，即土星的希腊名称）的缩写\(^\text{[29]}\)。该符号后演化出类似小写希腊字母 eta 的形状，并在 16 世纪顶部添加十字，以使这一原本的异教符号基督化。

罗马人以土星之名将一周的第七天命名为 Saturday，即 Sāturni diēs，“土星之日”\(^\text{[30]}\)。
\subsection{物理特征}
\begin{figure}[ht]
\centering
\includegraphics[width=8cm]{./figures/91b99dbcbbce021d.png}
\caption{土星与地球及地球月球大小的比较} \label{fig_Saturn_2}
\end{figure}
土星是一颗气态巨行星，主要由氢和氦组成。它缺乏明确的固体表面，但行星内部很可能存在一颗固体核心\(^\text{[31]}\)。由于自转的影响，土星呈现扁球体形状——两极被压扁、赤道区域鼓起。其赤道半径比极半径大超过 10\%：分别为 60,268 公里与 54,364 公里（37,449 英里与 33,780 英里）\(^\text{[6]}\)。

太阳系中的其他巨行星——木星、天王星与海王星——其扁率均低于土星。赤道凸起与自转速率的共同作用导致土星赤道处的有效表面重力仅为 8.96 m/s²，相当于其两极重力的 74\%，且低于地球表面重力。然而，赤道处的逃逸速度接近 36 km/s，远高于地球的逃逸速度\(^\text{[32]}\)。

土星是太阳系中唯一密度低于水的行星——约低 30\%\(^\text{[33]}\)。尽管其核心的密度远高于水，但由于其厚重大气层的存在，整颗行星的平均密度仅为 0.69 g/cm³。木星的质量是地球的 318 倍\(^\text{[34]}\)，而土星的质量是地球的 95 倍\(^\text{[6][35][36][37][38][39]}\)。木星与土星合计占据了太阳系全部行星质量的 92\%\(^\text{[40]}\)。
\subsubsection{内部结构}
\begin{figure}[ht]
\centering
\includegraphics[width=8cm]{./figures/0e870b2d45e0c034.png}
\caption{按比例绘制的土星示意图} \label{fig_Saturn_3}
\end{figure}
尽管土星主要由氢和氦组成，但其大部分质量并非处于气态，因为当密度超过 0.01 g/cm³ 时，氢会成为非理想液体，而这一密度在包含土星 99.9\% 质量的半径范围内即可达到。土星内部的温度、压力与密度在向核心方向不断上升，使得在更深层中氢呈金属态\(^\text{[40]}\)。

标准的行星模型表明，土星内部与木星相似，均包含一颗小型岩石核心，外围包裹着氢和氦，并含有少量各种挥发物\(^\text{[41]}\)。对土星形变的分析显示，其中心物质的集中程度显著高于木星，因此在其中心附近含有更多密度高于氢的物质。土星中心区域按质量计算约有 50\% 为氢，而木星约为 67\%\(^\text{[42]}\)。

该核心在组成上与地球类似，但密度更大。结合土星的引力矩测量结果与其内部物理模型，科学家得以对土星核心的质量设定约束。2004 年的估计认为其核心质量必须为地球质量的 9–22 倍\(^\text{[43][44]}\)，相当于直径约 25,000 公里（16,000 英里）\(^\text{[45]}\)。对土星光环的测量则显示其核心可能更加弥散，质量约等于 17 个地球，半径约为土星整体半径的 60\%\(^\text{[46]}\)。其外层包裹着更厚的液态金属氢层，然后是被氦饱和的液态分子氢层，随高度升高逐渐过渡为气体。最外层约有 1,000 公里（620 英里）厚，由气体组成\(^\text{[47][48][49]}\)。

土星内部非常炽热，其核心温度可达 11,700 °C（21,100 °F），并向外辐射的能量是其自太阳接收能量的 2.5 倍。木星的热能由缓慢的引力压缩（开尔文–亥姆霍兹机制）产生；但这一机制单独可能不足以解释土星的热量来源，因为土星质量更小。另一种或额外的机制可能是土星深层内部的氦滴“降雨”过程。当氦滴向密度较低的氢层下降时，摩擦会释放能量，并使土星外层的氦含量减少\(^\text{[50][51]}\)。这些下降的氦滴可能最终形成包围核心的氦层\(^\text{[41]}\)。在土星内部，如同在木星\(^\text{[52]}\)、以及冰巨行星天王星和海王星\(^\text{[53]}\)中一样，可能会出现钻石雨。
\subsubsection{大气层}  
土星外层大气按体积分数包含 96.3\% 的分子氢和 3.25\% 的氦。与太阳中该元素的丰度相比，土星大气中的氦比例显著偏低\(^\text{[41]}\)。比氦更重的元素（金属丰度）的精确含量尚不明确，但通常认为其比例与太阳系形成时期的原始丰度一致。估计这些较重元素的总质量为地球质量的 19–31 倍，其中相当一部分位于土星的核心区域\(^\text{[54]}\)。

在土星大气中已探测到微量的氨、乙炔、乙烷、丙烷、磷化氢和甲烷\(^\text{[55][56][57]}\)。上层云由氨晶体组成，而下层云似乎由硫氢化铵（NH\(_4\)SH）或水构成\(^\text{[58]}\)。太阳的紫外辐射会在上层大气中引发甲烷的光解，导致一系列烃类化学反应，其生成物再通过湍流与扩散向下输送。此光化学循环受土星年度季节周期的调制\(^\text{[57]}\)。卡西尼号曾观测到一系列出现在北纬的云层结构，被昵称为“珍珠串”。这些结构是位于更深云层中的云隙\(^\text{[59]}\)。

\textbf{云层}
\begin{figure}[ht]
\centering
\includegraphics[width=8cm]{./figures/5b480d0d41a65c33.png}
\caption{2011 年，一场全球性的风暴环绕整个行星。该风暴沿着行星一周移动，使得风暴的头部（明亮区域）绕行并经过其尾部。} \label{fig_Saturn_4}
\end{figure}
土星的大气呈现出类似木星的带状结构，但土星的云带要更为暗淡，并且在赤道附近宽度更大。用于描述这些云带的命名法与木星相同。直到 20 世纪 80 年代“旅行者”号飞掠任务之后，人们才观察到土星更细致的云图案。自那以后，地基望远镜观测技术已有显著提升，使得对土星的定期观测成为可能。(^\text{[60]})

云层的组成随深度和压力的增加而变化。在上层云层中，温度位于 100–160 K 范围内，压力约为 0.5–2 bar，云层由氨冰组成。水冰云层开始于约 2.5 bar 的压力位置，并向下延伸至 9.5 bar，此处的温度范围为 185–270 K。在这一层中混杂有一条硫氢化铵冰带，处于 3–6 bar 的压力范围内，温度为 190–235 K。更低的位置，压力在 10–20 bar、温度 270–330 K 的区域，则含有由水滴和溶于水中的氨构成的混合层。(^\text{[61]})

通常较为平淡的土星大气偶尔也会呈现出类似木星的长期存在的椭圆云以及其他特征。1990 年，“哈勃”太空望远镜拍摄到赤道附近一片巨大的白色云团，而在“旅行者”号飞掠期间这一特征并不存在；1994 年又观测到另一场较小的风暴。1990 年的风暴是“大白斑”的一个典型例子，这是一种短暂现象，约每土星年（约合地球 30 年）在北半球夏至前后出现一次。(^\text{[62]})

早期的“大白斑”曾分别在 1876 年、1903 年、1933 年和 1960 年被观测到，其中 1933 年的风暴观测最为详细。(^\text{[63]}) 最近一次大型风暴出现在 2010 年。2015 年，研究者利用甚大阵列（VLA）望远镜研究了土星大气，报告称他们发现了“所有中纬度大型风暴的长存痕迹、存在数百年历史的赤道风暴混合体，以及可能在北纬 70° 处存在一场未被记录的更古老的风暴”。(^\text{[64]})

土星的风速是太阳系行星中第二快的，仅次于海王星。“旅行者”号的数据表明其东风峰值风速可达 500 m/s（1,800 km/h）。(^\text{[65]}) 在“卡西尼”号于 2007 年拍摄的图像中，土星北半球呈现出类似天王星的亮蓝色，这种颜色很可能由瑞利散射造成。(^\text{[66]}) 热成像显示土星南极存在一个温暖的极涡，这是太阳系中唯一已知的此类现象。(^\text{[67]}) 虽然土星的典型温度约为 −185 °C，但极涡中的温度经常可达 −122 °C，被认为是土星上最温暖的区域。(^\text{[67]})

\textbf{六边形云层结构}
\begin{figure}[ht]
\centering
\includegraphics[width=8cm]{./figures/77d74f1601865d59.png}
\caption{土星南北两极的红外图像} \label{fig_Saturn_5}
\end{figure}
在约 78°N 纬度处，围绕北极涡旋的大气中存在一个持久的六边形波形结构，最早在旅行者号的图像中被发现\(^\text{[68][69][70]}\)。六边形的每一条边约为 14{,}500\,\text{km}（9{,}000\,\text{mi}），比地球的直径还要长\(^\text{[71]}\)。整个结构以 10\,\text{h}\,39\,\text{m}\,24\,\text{s} 的周期旋转（与该行星的无线电辐射周期一致），这一周期被认为等同于土星内部的自转周期\(^\text{[72]}\)。与土星大气中其他可见云层不同，这一六边形结构在经度方向上并不发生漂移\(^\text{[73]}\)。这一图案的成因仍是高度推测的话题。多数科学家认为它是大气中的驻波结构。通过差速旋转流体的实验，可以在实验室中复制多边形形状\(^\text{[74][75]}\)。

哈勃太空望远镜对南极区域的成像表明，该区域存在一条急流，但并未发现强极涡或任何六边形驻波\(^\text{[76]}\)。NASA 于 2006 年 11 月报道，卡西尼号观测到一个“类飓风”风暴，其中心固定在南极，并具有一个清晰的风眼墙\(^\text{[77][78]}\)。在地球之外的任何其他行星上此前均未观测到风眼墙。例如，加利略号拍摄的木星大红斑图像中并未显示风眼墙\(^\text{[79]}\)。

南极风暴可能已经存在了数十亿年\(^\text{[80]}\)。该极涡的规模可与地球相当，风速可达 550\,\text{km/h}\(^\text{[80]}\)。
\subsubsection{磁层}
\begin{figure}[ht]
\centering
\includegraphics[width=8cm]{./figures/d9afeaaa48604fc3.png}
\caption{土星北极的极光} \label{fig_Saturn_6}
\end{figure}
土星具有一个固有的磁场，其形状简单且对称——即一个磁偶极子。其在赤道处的强度为 0.2 高斯（20~μT），大约是木星磁场强度的二十分之一，并且比地球的磁场稍弱\(^\text{[27]}\)。因此，土星的磁层远小于木星的磁层\(^\text{[81]}\)。

当旅行者 2 号进入磁层时，太阳风压较高，磁层仅延伸到 19 个土星半径处，即 110 万千米（684,000 英里）\(^\text{[82]}\)，但在几个小时内磁层扩张，并保持这种状态约三天\(^\text{[83]}\)。磁场很可能与木星的磁场生成方式相似——由液态金属氢层中的电流（称为金属氢发电机）所产生\(^\text{[81]}\)。该磁层在偏转来自太阳的太阳风粒子方面十分有效。土星的卫星泰坦在土星磁层的外部区域轨道运行，并通过其外层大气中的电离粒子向磁层补充等离子体\(^\text{[27]}\)。与地球类似，土星的磁层也会产生极光\(^\text{[84]}\)。
\subsection{轨道与自转}
\begin{figure}[ht]
\centering
\includegraphics[width=8cm]{./figures/4ee15093919e1056.png}
\caption{土星及太阳系外行星绕太阳运行的动画} \label{fig_Saturn_7}
\end{figure}
\begin{figure}[ht]
\centering
\includegraphics[width=8cm]{./figures/23c0c3dad6383630.png}
\caption{土星在一次轨道运行周期中（2001–2029 年），从地球（在冲日位置）观测到的模拟外观} \label{fig_Saturn_8}
\end{figure}
土星与太阳之间的平均距离超过 14 亿千米（9 AU）。其平均轨道速度为 9.68 km/s，\(^\text{[6]}\) 土星需要 10,759 个地球日（约 29\(\tfrac{1}{2}\) 年）\(^\text{[85]}\) 才能完成一次绕太阳的公转。\(^\text{[6]}\) 因此，它与木星形成近似的 5:2 平均运动共振。\(^\text{[86]}\) 土星的椭圆轨道相对于地球轨道平面倾斜 2.48°。\(^\text{[6]}\) 其近日点与远日点距离平均分别为 9.195 与 9.957 AU。\(^\text{[6][87]}\) 土星上可见的大气特征会随纬度不同以不同速率自转，因此与木星类似，不同区域被分配了不同的自转周期。

天文学家使用三套不同的系统来描述土星的自转速率。系统 I 的自转周期为 10h 14m 00s（844.3°/d），包括赤道带、南赤道带与北赤道带。极区自转速率被认为与系统 I 类似。除南北极区以外的其他土星纬度称为系统 II，其自转周期被指定为 10h 38m 25.4s（810.76°/d）。系统 III 则指土星内部的自转速率。根据旅行者 1 号与旅行者 2 号探测到的土星无线电辐射，\(^\text{[88]}\) 系统 III 的自转周期为 10h 39m 22.4s（810.8°/d）。系统 III 已在很大程度上取代了系统 II。\(^\text{[89]}\)

土星内部自转周期的精确数值仍难以确定。卡西尼号于 2004 年接近土星时发现土星无线电自转周期明显增加，约为 10h 45m 45s ± 36s。\(^\text{[90][91]}\) 结合卡西尼号、旅行者号与先驱者号的多项测量结果，对土星整体自转速率的估计为 10h 32m 35s。\(^\text{[92]}\) 对土星 C 环的研究得出的自转周期为 10h 33m 38s \(^{+1m 52s}_{-1m 19s}\)。\(^\text{[17][18]}\)

2007 年 3 月，人们发现土星无线电辐射的变化与其自转速率不符。这种差异可能由土卫二（Enceladus）上的喷泉活动引起。喷泉活动向土星轨道释放的水蒸气被电离后，会对土星磁场施加拖曳，使得磁场自转相对于行星本体自转略微减慢。\(^\text{[93][94][95]}\)

土星目前仅已知一颗特洛伊小行星，即 2019 UO\(_{14}\)，其特洛伊构型于 2024 年 9 月公布，围绕太阳运行于土星轨道前方 60° 的稳定 L\(_4\) 拉格朗日点。\(^\text{[96]}\) 这一发现使得水星成为唯一尚未发现特洛伊天体的行星。轨道共振机制，包括岁差共振，被认为是已知土星特洛伊天体数量极少的原因之一。\(^\text{[97]}\)
\subsection{天然卫星}
\begin{figure}[ht]
\centering
\includegraphics[width=10cm]{./figures/72eb16859e562019.png}
\caption{艺术家笔下的土星、其光环以及主要的冰质卫星——从米玛斯到瑞亚} \label{fig_Saturn_9}
\end{figure}
土星已知有 274 颗卫星，\(^\text{[98][99][100][101]}\) 其中 63 颗拥有正式名称。\(^\text{[12][11]}\) 有证据显示，在土星光环中存在直径 40–500 米的数十到数百个微小卫星体，\(^\text{[102]}\) 但它们不被视为真正的卫星。最大卫星——泰坦（Titan）——占据土星环绕系统（包括光环）总质量的 90\% 以上。\(^\text{[103]}\) 土星第二大卫星——瑞亚（Rhea）——可能拥有其自身的稀薄光环系统，\(^\text{[104]}\) 并具有极其稀薄的大气层。\(^\text{[105][106][107]}\)

其他许多卫星都很小：其中 131 颗的直径小于 50 公里。\(^\text{[108]}\) 传统上，土星的大多数卫星都以希腊神话中的泰坦命名。泰坦是太阳系中唯一拥有主要大气层的卫星，\(^\text{[109][110]}\) 其中存在复杂的有机化学反应。它也是唯一拥有烃类湖泊的卫星。\(^\text{[111][112]}\)

2013 年 6 月 6 日，IAA-CSIC 的科学家报告在泰坦上层大气中探测到多环芳香烃（PAHs），这可能是生命的前体物质。\(^\text{[113]}\) 2014 年 6 月 23 日，美国宇航局（NASA）宣布有力证据表明：泰坦大气层中的氮并非来自形成土星的原始物质，而是来自与彗星有关的奥尔特云物质。\(^\text{[114]}\)

土卫二（Enceladus）在化学成分上似乎与彗星类似，\(^\text{[115]}\) 因而长期被视为可能存在微生物生命的栖息地。\(^\text{[116][117][118][119]}\) 支持这一可能性的证据包括：其富含盐分的颗粒具有“类海洋”的化学组成，表明土卫二喷出的多数冰来自盐水液体蒸发。\(^\text{[120][121][122]}\) 卡西尼号在 2015 年穿越土卫二羽状喷流时，探测到维持产甲烷生命形式所需的主要成分。\(^\text{[123]}\)

2014 年 4 月，美国宇航局科学家报告称：土星 A 环中可能正在形成一颗新的卫星，此前卡西尼号在 2013 年 4 月 15 日的图像中拍摄到了该结构。\(^\text{[124]}\)
\subsection{行星环}
\begin{figure}[ht]
\centering
\includegraphics[width=8cm]{./figures/8c0274dc32d85fc0.png}
\caption{v土星的光环（此处为卡西尼号于 2004 年 10 月拍摄）是太阳系中最庞大、最显眼的光环系统。\(^\text{[48]}\)} \label{fig_Saturn_10}
\end{figure}
土星最为人所知的，或许就是其行星光环系统，使其在视觉上独具特色。\(^\text{[48]}\) 光环自土星赤道向外延伸，从 6,630 千米至 120,700 千米（4,120–75,000 英里），平均厚度约为 20 米（66 英尺）。光环主要由水冰组成，并含有微量的托林（tholin）杂质以及约 7\% 的无定形碳颗粒作为表面覆盖。\(^\text{[125]}\) 构成光环的颗粒大小从微小尘埃到直径 10 米不等。\(^\text{[126]}\) 虽然其他气体巨行星也拥有光环系统，但土星的光环最大且最为醒目。

关于光环的年龄存在争论。一种观点认为光环非常古老，与土星一起由原始星云物质形成（约 46 亿年前），\(^\text{[127]}\) 或形成于晚期重轰炸（LHB）之后不久（约 41–38 亿年前）。\(^\text{[128][129]}\) 另一种观点认为光环年轻得多，大约形成于 1 亿年前。\(^\text{[130][131][132]}\) 支持后者理论的麻省理工学院研究团队提出，光环是土星一颗被摧毁的卫星——名为“Chrysalis”——的残余物。\(^\text{[133]}\)

在主光环之外，在距离行星 1,200 万千米（750 万英里）处存在稀疏的菲比（Phoebe）光环。该光环相对于其他光环倾斜 27°，并像菲比一样以逆行方式运行。\(^\text{[134]}\)

土星的一些卫星，如潘多拉（Pandora）和普罗米修斯（Prometheus），作为“牧羊卫星”帮助限制光环、阻止其向外扩散。\(^\text{[135]}\) 潘（Pan）和阿特拉斯（Atlas）会在土星光环中激发弱的线性密度波，这些密度波已用于更可靠地计算它们的质量。\(^\text{[136]}\)
\begin{figure}[ht]
\centering
\includegraphics[width=14.25cm]{./figures/bf65e1a518a4adfe.png}
\caption{卡西尼号窄角相机于 2007 年 5 月 9 日拍摄的土星 D、C、B、A 与 F 环（从左至右）未受光照一侧的自然色拼接图像。图中距离均以土星中心为基准。} \label{fig_Saturn_11}
\end{figure}
\subsection{观测与探测的历史}
对土星的观测与探测可分为三个阶段：（1）前现代时期肉眼观测；（2）自 17 世纪起从地球进行的望远镜观测；（3）由空间探测器进行的访问，包括飞越与轨道探测。在 21 世纪，对土星的望远镜观测仍在继续（包括地球上的望远镜以及哈勃太空望远镜等地球轨道天文台），并由卡西尼号在土星轨道进行观测直至其于 2017 年退役。
\subsubsection{前望远镜时代的观测}  
参见：Saturn（神话）以及 占星术中的行星 § Saturn  
土星自史前时代以来就为人所知，\(^\text{[137]}\) 在早期的文字记录中，它是多种神话中的重要角色。巴比伦天文学家系统地观测并记录了土星的运动。\(^\text{[138]}\) 在古希腊语中，该行星被称为 Φαίνων Phainon，\(^\text{[139]}\) 而在罗马时代，它被称为“土星之星”或“太阳（即 Helios）之星”。\(^\text{[140][141]}\) 在古罗马神话中，Phainon 这颗行星被视为农业之神的神圣星体，现代名称即源于此。\(^\text{[142]}\) 罗马人认为土星神（Saturnus）与希腊神克罗诺斯（Cronus）相对应。在现代希腊语中，这颗行星仍保留克罗诺斯之名——Κρόνος：Kronos。\(^\text{[143]}\)

希腊科学家托勒密根据其在土星处于冲日位置时所做的观测，推算了土星的轨道。\(^\text{[144]}\) 在印度占星术中，有九大占星天体，被称为 Navagrahas。土星被称为 “Shani”，并根据人生中的善恶行为对每个人作出评判。\(^\text{[142][144]}\) 古代中国与日本文化将土星称为“土星”，源自用于分类自然元素的五行体系。\(^\text{[145][146][147]}\)

在希伯来语中，土星被称为 Shabbathai。\(^\text{[148]}\) 其守护天使为 Cassiel。其智慧体或善灵为 ‘Agȋȇl（希伯来语：אגיאל，拉丁化：ʿAgyal），\(^\text{[149]}\) 而其阴暗灵（恶魔）为 Zȃzȇl（希伯来语：זאזל，拉丁化：Zazl）。\(^\text{[149][150][151]}\) Zazel 被描述为在所罗门魔法中被召唤的伟大天使，“在爱情召唤中十分有效”。\(^\text{[152][153]}\) 在奥斯曼土耳其语、乌尔都语与马来语中，Zazel 的名字为 ‘Zuhal’，源自阿拉伯语（阿拉伯语：زحل，拉丁化：Zuhal）。\(^\text{[150]}\)
\subsubsection{望远镜时代的太空飞行前观测}
\begin{figure}[ht]
\centering
\includegraphics[width=8cm]{./figures/de1ff6ea3f6b74db.png}
\caption{罗伯特·胡克在 1666 年的这幅土星绘图中注意到球体与光环彼此投下的阴影（a 与 b）。} \label{fig_Saturn_12}
\end{figure}
\begin{figure}[ht]
\centering
\includegraphics[width=8cm]{./figures/b8d92d1f5c26ae9a.png}
\caption{伽利略·伽利莱在 1610 年观测到了土星的光环，但当时无法判断它们的本质。} \label{fig_Saturn_13}
\end{figure}
土星的光环至少需要一架口径 15 毫米的望远镜才能分辨，\(^\text{[154]}\) 在此之前它们并不为人所知，直到克里斯蒂安·惠更斯于 1655 年首次看到，并于 1659 年发表观测结果。伽利略在 1610 年用他原始的望远镜观测土星时，\(^\text{[155][156]}\) 错误地将土星看起来不像圆形的现象解释为其两侧各有一颗卫星。\(^\text{[157][158]}\)

当惠更斯使用更高的望远镜放大倍率时，这种看法被推翻，人类首次真正看清了土星光环。惠更斯还发现了土星的卫星——泰坦。后来，焦万尼·多梅尼科·卡西尼又发现了另外四颗卫星：土卫八（Iapetus）、土卫五（Rhea）、土卫三（Tethys）和土卫四（Dione）。1675 年，卡西尼发现了如今称为“卡西尼缝隙”的光环间隙。\(^\text{[159]}\)

此后直到 1789 年，没有进一步的重要发现，直到威廉·赫歇尔发现了另外两颗卫星：米玛斯（Mimas）和恩克拉多斯（Enceladus）。1848 年，一支英国团队发现了不规则形状的卫星海帕里翁（Hyperion），该卫星与泰坦存在轨道共振关系。\(^\text{[160]}\)

1899 年，威廉·亨利·皮克林发现了菲比（Phoebe），这是一颗高度不规则的卫星，与土星的大型卫星不同，它不是同步自转。\(^\text{[160]}\) 菲比是首颗被发现的此类卫星，它以逆行轨道绕土星运行，需要一年多才能完成一次环绕。20 世纪初，对泰坦的研究在 1944 年确认了它拥有厚重大气——这是太阳系众卫星中独一无二的特征。\(^\text{[161]}\)
\subsubsection{太空飞行任务}  
\textbf{先驱者 11 号飞掠}
\begin{figure}[ht]
\centering
\includegraphics[width=8cm]{./figures/b0f551c206af9e74.png}
\caption{先驱者 11 号拍摄的土星图像} \label{fig_Saturn_14}
\end{figure}
先驱者 11 号于 1979 年 9 月首次飞掠土星，当时其距离行星云顶仅 20,000 千米（12,000 英里）。飞船拍摄了土星及其几颗卫星的图像，但分辨率过低，无法辨认其表面细节。飞船还研究了土星光环，揭示了狭窄的 F 环，并发现光环中的暗缝隙在高相角（朝向太阳）观测时会显得明亮，这意味着其中含有能散射光的细小颗粒。此外，先驱者 11 号还测量了泰坦的温度。\(^\text{[162]}\)

旅行者号飞掠  
1980 年 11 月，旅行者 1 号探测器访问了土星系统。它传回了首批高分辨率的土星、其光环及其卫星的图像。多颗卫星的表面特征由此首次被看见。旅行者 1 号对泰坦进行了近距离飞掠，使人们对这颗卫星的大气层有了更深入的了解。它证实泰坦的大气在可见光波段完全不透明；因此未能看到其表面细节。此次飞掠使探测器的轨道偏离了太阳系的主平面。\(^\text{[163]}\)

将近一年后的 1981 年 8 月，旅行者 2 号继续研究土星系统。获得了更多土星卫星的特写图像，并观测到大气与光环发生变化的证据。在飞掠期间，探测器的可转动摄像平台曾卡住数日，导致部分计划中的成像任务丢失。土星的引力被用来将探测器引向天王星。\(^\text{[163]}\)

旅行者号探测器发现并确认了数颗在土星光环附近或其中运行的新卫星，并首次识别了麦克斯韦间隙（C 环中的一个间隙）以及凯勒间隙（A 环中一个宽 42 千米的间隙）。\(^\text{[164]}\)

\textbf{卡西尼–惠更斯号飞船}
\begin{figure}[ht]
\centering
\includegraphics[width=8cm]{./figures/ea4c046d929e30e7.png}
\caption{在土卫二的南极，喷泉状的间歇泉沿着“虎纹”裂缝的多处位置喷射出水。\(^\text{[165]}\)} \label{fig_Saturn_15}
\end{figure}
卡西尼–惠更斯号太空探测器于 2004 年 7 月 1 日进入土星轨道。2004 年 6 月，它对菲比进行了近距离飞掠，并传回高分辨率图像与数据。卡西尼在飞掠土星最大卫星泰坦时，利用雷达拍摄到了大型湖泊及其海岸线，并观测到众多岛屿与山脉。轨道器在释放惠更斯探测器之前完成了两次泰坦飞掠，惠更斯于 2004 年 12 月 25 日释放，并于 2005 年 1 月 14 日降落在泰坦表面。\(^\text{[166]}\)

自 2005 年初起，科学家使用卡西尼号追踪土星上的闪电。其闪电能量约为地球闪电的 1,000 倍。\(^\text{[167]}\)

2006 年，美国宇航局（NASA）报告称，卡西尼号在土卫二上发现了距离地表不超过数十米的液态水储库，其喷发形成了喷泉状的间歇泉。这些冰质颗粒喷流由土卫二南极区域的喷口喷出，并进入环绕土星的轨道。\(^\text{[168]}\) 在土卫二上已识别出超过 100 处间歇泉。\(^\text{[165]}\) 2011 年 5 月，NASA 科学家报告称，按照人类目前对生命的理解，土卫二“正逐渐显现为地球之外太阳系中最适宜生命存在的地点”。\(^\text{[169][170]}\)
\begin{figure}[ht]
\centering
\includegraphics[width=8cm]{./figures/d10622179066806d.png}
\caption{从卡西尼号视角看到的土星掩日景象。光环清晰可见，包括 F 环。} \label{fig_Saturn_16}
\end{figure}
卡西尼号的照片揭示了一个此前未被发现的行星光环，位于土星明亮主光环之外、G 环与 E 环之间。该光环的来源被推测为一颗流星体撞击雅努斯（Janus）或艾比米修斯（Epimetheus）后所产生的物质。\(^\text{[171]}\)

2006 年 7 月，卡西尼号传回了泰坦北极附近烃类湖泊的影像，并在 2007 年 1 月得到确认。2007 年 3 月，在泰坦北极附近发现了烃类海洋，其中最大的海洋面积几乎与里海相当。\(^\text{[172]}\)

2006 年 10 月，探测器在土星南极探测到一场直径达 8,000 千米（5,000 英里）的类气旋风暴，其结构具有“风眼壁”（eyewall）。\(^\text{[173]}\)

自 2004 年至 2009 年 11 月 2 日，该探测器共发现并确认了八颗新的土星卫星。\(^\text{[174]}\) 2013 年 4 月，卡西尼号传回了土星北极风暴的影像，其规模比地球上的飓风大 20 倍，风速超过 530 km/h（330 mph）。\(^\text{[175]}\)

2017 年 9 月 15 日，卡西尼–惠更斯号探测器执行了其任务的“壮丽终章”（Grand Finale）：多次穿越土星与内环之间的狭窄缝隙。\(^\text{[176][177]}\) 卡西尼号的最终大气入射（atmospheric entry）宣告了任务的结束。

\textbf{可能的未来任务}

继续探索土星仍被 NASA 视为可行选择，并可能纳入其进行中的“新前沿”（New Frontiers）任务计划。NASA 先前曾征求包含“土星大气进入探测器”（Saturn Atmospheric Entry Probe）的任务方案，并考虑通过“蜻蜓号”（Dragonfly）探测泰坦与土卫二（Enceladus）的宜居性与可能存在生命的证据。\(^\text{[178][179]}\)
\subsection{观测}
\begin{figure}[ht]
\centering
\includegraphics[width=8cm]{./figures/6e9ae23ff5c5de0e.png}
\caption{业余望远镜看到的土星影像} \label{fig_Saturn_17}
\end{figure}
土星是从地球肉眼可轻易看到的五颗行星中距离最远的一颗，其余四颗为水星、金星、火星与木星。（天王星以及在暗夜条件下偶尔可见的 4 号灶神星也能被肉眼看到。）土星在夜空中呈现为一颗明亮的、略带黄色的光点。土星的平均视星等为 0.46，标准差为 0.34。\(^\text{[24]}\) 其亮度的大部分变化来源于光环相对于太阳与地球的倾角变化。当光环倾角最大时，其最亮视星等为 −0.55；而当光环几乎与视线平行时，其最暗视星等为 1.17。\(^\text{[24]}\)

土星在黄道背景星座中完成一次完整循环需要约 29.4 年。大多数人需要至少放大 30 倍的光学仪器（大型双筒望远镜或小型望远镜）才能分辨出土星光环的清晰影像。\(^\text{[48][154]}\)

当地球穿过光环平面时（每个土星年发生两次，大约每 15 个地球年），由于光环极薄，它们会在短暂时间内从视野中“消失”。最近的一次“消失”发生在 2025 年，但当时土星过于接近太阳，无法观测。\(^\text{[180]}\)

土星及其光环在行星处于冲日（即延角为 180°、在天空中与太阳相对）或其附近位置时最易观测。土星的冲日每年发生一次——约每 378 天——此时土星最为明亮。由于地球与土星都绕太阳运行在离心轨道上，它们与太阳的距离随时间变化，因此彼此之间的距离也随之变化，导致一次冲日与下一次冲日之间土星亮度有所不同。当光环倾角更大、更加显露时，土星也会显得更亮。例如，在 2002 年 12 月 17 日的冲日期间，由于光环的有利朝向，土星呈现出其最亮状态，\(^\text{[181]}\) 尽管于 2003 年末土星距离地球与太阳都更近。\(^\text{[181]}\)

土星偶尔会被月球掩食（即月球在天空中遮挡土星）。与太阳系所有行星一样，土星掩食呈“季节性”发生。一个掩食周期中，大约会有 12 个月的月度掩食，随后约五年不会再发生掩食。由于月球轨道相对于土星轨道倾斜数度，只有当土星接近两轨道交点时才会发生掩食（这一周期同时受土星轨道周期与月球轨道的 18.6 地球年交点岁差周期所控制）。\(^\text{[182]}\)
\subsection{小说中的土星}   
自至少 1752 年伏尔泰出版小说《Micromégas》以来，土星便频繁出现在文学作品中。\(^\text{[183]}\) 早期作品通常将其描绘为固体星球，\(^\text{[184]}\) 而后期作品则正确地将其描述为气态行星。土星的卫星也常出现在文学作品中，尤其是泰坦。\(^\text{[185]}\)
\subsection{参见} 
\begin{itemize}
\item 土星的卫星  
\item 太阳系行星统计  
\item 土星纲要
\end{itemize}  
\subsection{注释}\\ 
a.左下角的点是泰坦\\  
b.指 1 bar 大气压水平位置\\  
c.基于 1 bar 大气压水平内的体积计算
\subsection{参考资料}
\begin{enumerate}
\item 伊丽莎白·沃尔特 (2003年4月21日)。剑桥高级学习词典(第二版)。剑桥大学出版社。ISBN 978-0-521-53106-1。
\item "土星"。牛津英语词典(在线版)。牛津大学出版社。 (订阅或参与机构成员必需。)
\item “利用小型放射性同位素动力系统进行勘探”(PDF)。NASA.2004年9月。存档自原来的(PDF)2016年12月22日。已检索1月26日2016。
\item M ü ller等人 (2010年)。“克朗磁层中的方位等离子体流”。地球物理研究杂志。115(A8) 2009JA015122: a08203。Bibcode:2010JGRA..115.8203M。doi:10.1029/2009ja015122。ISSN0148-0227。 
\item "Cronian"。牛津英语词典(在线版)。牛津大学出版社。 (订阅或参与机构成员必需。)
\item 威廉姆斯，大卫R.(2016年12月23日)。"土星概况"。NASA.存档自原来的2017年7月17日。已检索10月12日2017。
\item “行星物理参数”。NASA喷气推进实验室。
\item 西蒙，J.L.；布列塔尼翁，P.；查普朗；查普朗-图兹，M.；弗朗库，G.；拉斯卡，J。(1994年2月)。“月球和行星的进动公式和平均元素的数值表达式”。天文学和天体物理学。282(2):663-683。Bibcode:1994A & A。282 .. 663S。
\item 苏阿米，D.；Souchay, J.(2012年7月)。“太阳系的不变平面”。天文学与天体物理学。543: 11。Bibcode:2012A & A...543A.133S。doi:10.1051/0004-6361/201219011。A133.
\item “2032年11月近日点的地平线行星中心批量呼叫”。ssd.jpl.nasa.gov(土星行星中心 (699) 的近日点发生在2032年11月29日，在rdot从负到正的翻转过程中在9.0149170au)。NASA/JPL。已存档从原始的2021年9月7日。已检索9月7日2021。
\item “太阳系动力学-行星卫星发现情况”。NASA.2021年11月15日。已检索6月4日2022年。
\item “土星现在以62颗新发现的卫星领先月球竞赛”。UBC科学。不列颠哥伦比亚大学。2023年5月11日。已检索5月11日2023。
\item “通过数字-土星”。NASA太阳系探索。NASA。2017年11月10日。已存档从2018年5月10日开始。已检索8月5日2020。
\item “NASA: 太阳系探索: 行星: 土星: 事实与数据”。Solarsystem.nasa.gov。2011年3月22日。存档自原来的2011年9月2日。已检索8月8日2011年。
\item 福特尼，j.j.；赫尔，R；奈特莱曼，N.；史蒂文森，d.j.；马利，m.s.; 哈伯德，W.B.；Iess, L.(2018年12月6日)。“土星的内部”。在贝恩斯，K.H.；弗拉萨尔，f.m.；克虏伯，N.；Stallard，T. (编辑)。21世纪的土星。剑桥大学出版社。pp. 44-68.ISBN 978-1-108-68393-7。已存档从原件到2020年5月2日。已检索7月23日2019。
\item 塞利格曼，考特尼。“轮换周期和日长”。已存档从2011年7月28日起。已检索8月13日2009年。
\item 麦卡特尼，格雷琴; 温德尔，乔安娜 (2019年1月18日)。“科学家终于知道土星上的时间了”。NASA。已检索1月3日2025。
\item 曼科维奇，克里斯托弗; 等人 (2019年1月17日)。“卡西尼环地震学作为土星内部的探测器。I.刚性旋转”。天体物理学杂志。871(1): 1。arXiv:1805.10286。Bibcode:2019年4月。871。1米。doi:10.3847/1538-4357/aaf798。S2CID67840660。 
\item 阿基纳，B。A.;阿克顿，C。H、; A'Hearn，M。F.;康拉德，A.；Consolmagno, G.J.;达克斯伯里，T.；赫斯特罗弗，D.；希尔顿，J。L.;柯克，R.L.;Klioner, S.A.;麦卡锡，D.；米奇，K.；奥伯斯特，J.；平，J.；Seidelmann，P。K。(2018)。“IAU制图坐标和旋转要素工作组的报告: 2015年”。天体力学与动力天文学。130(3): 22。Bibcode:2018 CeMDA.130. .. 22A。doi:10.1007/s10569-017-9805-5。ISSN0923-2958。 
\item Hanel，r.a.; 等人 (1983)。“土星的反照率、内部热通量和能量平衡”。伊卡洛斯。53(2):262-285。Bibcode:1983Icar...53 .. 262小时。doi:10.1016/0019-1035(83)90147-1。
\item 马拉玛，安东尼; 克鲁布塞克，布鲁斯; 巴夫洛夫，赫里斯托 (2017)。“行星的综合宽带震级和反照率，适用于系外行星和第九行星”。伊卡洛斯。282:19-33。arXiv:1609.05048。Bibcode:2017Icar .. 282...19米。doi:10.1016/j.icarus.2016.09.023。S2CID119307693。 
\item “土星的温度范围”。科学。2018年4月20日。已存档从2021年5月26日开始。已检索5月26日2021。
\item “土星”。国家气象局。已存档从2021年5月26日开始。已检索5月26日2021。
\item Mallama, A.;希尔顿，J.L.(2018)。“计算天文历书的视在行星星等”。天文学与计算。25:10-24.arXiv:1808.01973。Bibcode:2018A & C .... 25...10米。doi:10.1016/j.ascom.2018.08.002。S2CID69912809。 
\item “百科全书-最聪明的身体”。IMCCE。已检索5月29日2023。
\item Knecht，Robin (2005年10月24日)。“在不同行星的大气层上”(PDF)。存档自原来的(PDF)2017年10月14日。已检索10月14日2017。
\item 罗素，C。等 (1997年)。“土星: 磁场和磁层”。科学。207(4429):407-10。Bibcode:1980Sci。207。407s。doi:10.1126/science.207.4429.407。PMID17833549。S2CID41621423。已存档从2011年9月27日的原件。已检索4月29日2007。  
\item Munsell，Kirk (2005年4月6日)。“土星的故事”。美国宇航局喷气推进实验室; 加州理工学院。存档自原来的2008年8月16日。已检索7月7日2007。
\item 亚历山大·琼斯 (1999)。来自Oxyrhynchus的天文纸莎草。美国哲学学会.pp. 62-63.ISBN 978-0-87169-233-7。已存档从原件到2021年4月30日。已检索9月28日2021。
\item 迈克尔·福尔克 (1999年6月)，“一周中几天的天文名称”,加拿大皇家天文学会杂志,93:122-133,Bibcode:1999JRASC..93..122F,已存档从原件到2021年2月25日,检索到11月18日2020
\item Melosh，H. Jay (2011)。行星表面过程。剑桥行星科学。卷13.剑桥大学出版社。第5页。ISBN 978-0-521-51418-7。已存档从原始的2017年2月16日。已检索2月9日2016。
\item Gregersen，Erik，编辑 (2010)。外太阳系: 木星、土星、天王星、海王星和矮行星。罗森出版集团。第119页。ISBN 978-1-61530-014-3。已存档从原件到2020年6月10日。已检索2月17日2018。
\item “土星-我们太阳系最美丽的星球”。保留文章。2011年1月23日已存档从2012年1月20日开始。已检索7月24日2011年。
\item 威廉姆斯，大卫R.(2004年11月16日)。“木星概况介绍”。NASA.存档自原来的2011年9月26日。已检索8月2日2007。
\item Brainerd，Jerome James (2004年10月6日)。“太阳系行星与地球相比”。天体物理学的旁观者。存档自原来的2011年10月1日。已检索7月5日2010。
\item Dunbar，Brian (2007年11月29日)。“NASA-土星”。NASA.存档自原来的2011年9月29日。已检索7月21日2011年。
\item 该隐，弗雷泽 (2008年7月3日)。“土星的质量”。今天的宇宙。已存档从2011年9月26日的原件。已检索8月17日2011年。
\item 杰罗姆·詹姆斯·布雷纳德 (2004年11月24日)。“土星的特征”。天体物理学的旁观者。存档自原来的2011年10月1日。已检索7月5日2010。
\item "关于土星的一般信息"。Scienceray。2011年7月28日存档自原来的2011年10月7日。已检索8月17日2011年。
\item 乔纳森·J·福特尼；Nettelmann，Nadine (2010年5月)。“巨型行星的内部结构、组成和演化”。空间科学评论。152(1-4):423-447.arXiv:0912.0533。Bibcode:2010SSRv .. 152 .. 423F。doi:10.1007/s11214-009-9582-x。S2CID49570672。 
\item 吉洛，特里斯坦; 等人 (2009年)。“土星在卡西尼-惠更斯之外的探索”。在米歇尔·K·多尔蒂；拉里·W·埃斯波西托；Krimigis，Stamatios M. (编辑)。来自卡西尼-惠更斯的土星。施普林格科学 + 商业媒体b.V.第745页。arXiv:0912.2020。Bibcode:2009sfch.book .. 745G。doi:10.1007/978-1-4020-9217-6_23。ISBN 978-1-4020-9216-9。S2CID 37928810。
\item “土星-内部 | Britannica”。www.britannica.com。已检索4月14日2022年。
\item 福特尼，乔纳森·J.(2004)。“看着巨大的行星”。科学。305(5689):1414-1415。doi:10.1126/science.1101352。PMID15353790。S2CID26353405。已存档从原始的2019年7月27日。已检索6月28日2019。  
\item 索门，D.；内疚，T。(2004年7月)。“氘的冲击压缩以及木星和土星的内部”。天体物理学杂志。609(2):1170-1180。arXiv:astro-ph/0403393。Bibcode:2004年4月。609.1170s。doi:10.1086/421257。S2CID119325899。 
\item "土星"。英国广播公司。2000。已存档从2011年1月1日起。已检索7月19日2011年。
\item 克里斯托弗·R·曼科维奇；吉姆·富勒 (2021)。“环形地震学揭示的土星弥漫核心”。自然天文学。5(11):1103-1109。arXiv:2104.13385。Bibcode:2021NatAs...5.1103M。doi:10.1038/s41550-021-01448-3。S2CID233423431。已存档从2021年8月20日开始。已检索8月22日2021。 
\item 福雷，冈特; 梅辛，特蕾莎·M。(2007)。行星科学导论: 地质学视角。斯普林格。第337页。ISBN 978-1-4020-5233-0。已存档从原始的2017年2月16日。已检索2月9日2016。
\item "土星"。国家海事博物馆。2015年8月20日存档自原来的2008年6月23日。已检索7月6日2007。
\item “土星内部的结构”。宇宙的窗户存档自原来的2011年9月17日。已检索7月19日2011年。
\item `德佩特，Imke; 利绍尔，杰克J。(2010)。行星科学(第二版)。剑桥大学出版社。pp. 254-255.ISBN 978-0-521-85371-2。已存档从2017年2月17日的原件。已检索2月9日2016。
\item “NASA-土星”。NASA.2004。存档自原来的2010年12月29日。已检索7月27日2007。
\item 克莱默，米里亚姆 (2013年10月9日)。钻石雨可能会填满木星和土星的天空。Space.com。已存档从原始的2017年8月27日。已检索8月27日2017。
\item 莎拉·卡普兰 (2017年8月25日)。“它在天王星和海王星上下雨固体钻石”。华盛顿邮报。已存档从原始的2017年8月27日。已检索8月27日2017。
\item 特里斯坦·吉洛 (1999)。“太阳系内外巨型行星的内部”。科学。286(5437):72-77.Bibcode:1999Sci...286...72G。doi:10.1126/science.286.5437.72。PMID10506563。S2CID6907359。  
\item Courtin，R.; 等人 (1967)。“从旅行者虹膜光谱看北纬温带土星大气的组成”。美国天文学会公报。15: 831。Bibcode:1983BAAS...15..831C。
\item 该隐，弗雷泽 (2009年1月22日)。“土星的气氛”。今天的宇宙。已存档从2012年1月12日起。已检索7月20日2011年。
\item Guerlet, S.;福切特，T.；Bézard, B.(2008年11月)。夏邦内尔，C.；梳齿，F.；Samadi，R. (编辑)。“卡西尼/CIRS肢体观测中土星平流层中的乙烷、乙炔和丙烷分布”。SF2A-2008: 法国天文学和天体物理学学会年会论文集: 405。Bibcode:2008sf2a.conf .. 405G。
\item 卡罗来纳州马丁内斯 (2005年9月5日)。“卡西尼号发现土星的动态云运行深”。NASA.已存档从2011年11月8日的原件。已检索4月29日2007。
\item “卡西尼号图像显示土星披着一串珍珠”(新闻稿)。卡罗来纳·马丁内斯，美国宇航局。2006年11月10日。已存档从2013年5月1日开始。已检索3月3日2013。
\item 奥顿，格伦·S.(2009年9月)。“为航天器探索外行星提供地面观测支持”。地球、月球和行星。105(2-4):143-152。Bibcode:2009EM & P。。105 .. 143O。doi:10.1007/s11038-009-9295-x。S2CID121930171。 
\item 米歇尔·K·多尔蒂；拉里·W·埃斯波西托；Krimigis, Stamatios M.(2009)。米歇尔·K·多尔蒂；拉里·W·埃斯波西托；Krimigis，Stamatios M. (编辑)。来自卡西尼-惠更斯的土星。斯普林格。第162页。Bibcode:2009sfch.book ..... D。doi:10.1007/978-1-4020-9217-6。ISBN 978-1-4020-9216-9。已存档从原始的2017年4月16日。已检索2月9日2016。
\item 佩雷斯-霍约斯，S.；桑切斯-拉韦格，A.；法语，R.G.;J.F.,罗哈斯 (2005)。“来自哈勃太空望远镜十年图像 (1994-2003) 的土星云结构和时间演变”。伊卡洛斯。176(1):155-174.Bibcode:2005年Icar。。176。。155便士。doi:10.1016/j.icarus.2005.01.014。
\item Kidger，Mark (1992)。“1990年的土星大白斑”。在帕特里克·摩尔(ed.)。1993年天文学年鉴。伦敦: W.诺顿公司。pp. 176-215.Bibcode:1992年ybas.conf ..... M。
\item 李，程; 德佩特，伊姆克; 莫克尔，克里斯; 索特，R。J.;巴特勒，布莱恩; 德波尔，大卫; 张志盟 (2023年8月11日)。“土星巨大风暴的持久、深刻影响”。科学进展。9(32) eadg9419。Bibcode:2023SciA .... 9G9419L。doi:10.1126/sciadv.adg9419。PMC10421028。PMID37566653。  
\item 汉密尔顿，卡尔文J.(1997)。“旅行者土星科学摘要”。Solarviews。存档自原来的2011年9月26日。已检索7月5日2007。
\item 渡边，苏珊 (2007年3月27日)。“土星的奇异六边形”。NASA.已存档从2010年1月16日的原件。已检索7月6日2007。
\item “土星上的温暖极地涡旋”。梅里尔维尔社区天文馆。2007。存档自原来的2011年9月21日。已检索7月25日2007。
\item 戈弗雷，D。A.(1988)。“土星北极周围的六边形特征”。伊卡洛斯。76(2): 335。Bibcode:1988年Icar。76 .. 335克。doi:10.1016/0019-1035(88)90075-9。
\item Sanchez-lavega，A.; 等人 (1993年)。“土星北极点和六边形的地面观测”。科学。260(5106):329-32。Bibcode:1993Sci。260 .. 329s。doi:10.1126/science.260.5106.329。PMID17838249。S2CID45574015。  
\item 再见，丹尼斯(2014年8月6日)。“风暴追逐土星”。纽约时报。已存档从原件到2018年7月12日。已检索8月6日2014。
\item “新图像显示了土星的怪异六边形云”。NBC新闻。2009年12月12日。已存档从原件到2020年10月21日。已检索9月29日2011年。
\item 戈弗雷，D。A.(1990年3月9日)。“土星极六边形的自转周期”。科学。247(4947):1206-1208.Bibcode:1990sci。247。1206克。doi:10.1126/science.247.4947.1206。PMID17809277。S2CID19965347。  
\item 贝恩斯，凯文·H; 等人 (2009年12月)。“卡西尼/维姆斯揭示的土星北极气旋和六角形深度”。行星与空间科学。57(14-15):1671-1681。Bibcode:2009P & SS。57.1671b。doi:10.1016/j.pss.2009.06.026。
\item 菲利普·鲍尔 (2006年5月19日)。“几何漩涡显露出来”。性质。doi:10.1038/news060515-17。S2CID129016856。 行星大气中漩涡中心出现的奇异几何形状可以用一桶水的简单实验来解释，但这与土星的模式相关联还不确定。
\item 阿吉亚尔，安娜·C。Barbosa; 等人 (2010年4月)。“土星北极六边形的实验室模型”。伊卡洛斯。206(2):755-763。Bibcode:2010Icar .. 206 .. 755B。doi:10.1016/j.icarus.2009.10.022。在液体溶液中旋转圆盘的实验室实验在类似于土星的稳定六边形图案周围形成涡流。
\item Sánchez-lavega，A.; 等人 (2002年10月8日)。“哈勃太空望远镜对1997年至2002年土星南极大气动力学的观测”。美国天文学会公报。34: 857。Bibcode:2002DPS....34.1307S。已存档从2010年6月30日的原件。已检索7月6日2007。
\item “图像PIA09187的NASA目录页面”。NASA行星摄影杂志。已存档从2011年11月9日的原件。已检索5月23日2007。
\item “巨大的 '飓风' 肆虐土星”。BBC新闻。2006年11月10日。已存档从2012年8月3日的原件。已检索9月29日2011年。
\item “NASA看到土星上的怪物风暴的眼睛”。NASA.2006年11月9日。存档自原来的2008年5月7日。已检索11月20日2006。
\item Nemiroff, R.;邦内尔，J.，编辑。(2006年11月13日)。“土星南极上空的飓风”。当天的天文图片。NASA。已检索5月1日2013。
\item 麦克德莫特，马修 (2000)。“土星: 大气层和磁层”。Thinkquest互联网挑战。已存档从2011年10月20日的原件。已检索7月15日2007。
\item “旅行者号-土星的磁层”。NASA喷气推进实验室2010年10月18日。存档自原来的2012年3月19日。已检索7月19日2011年。
\item 南希阿特金森 (2010年12月14日)。“热等离子体爆炸使土星的磁场膨胀”。今天的宇宙。已存档从2011年11月1日起。已检索8月24日2011年。
\item Russell，Randy (2003年6月3日)。“土星磁层概述”。宇宙的窗户存档自原来的2011年9月6日。已检索7月19日2011年。
\item 该隐，弗雷泽 (2009年1月26日)。“土星的轨道”。今天的宇宙。已存档从2011年1月23日起。已检索7月19日2011年。
\item 米奇琴科，T.A.;费拉兹-梅洛，S.(2001年2月)。“模拟木星-土星行星系统中的5: 2平均运动共振”。伊卡洛斯。149(2):357-374。Bibcode:2001Icar .. 149 .. 357米。doi:10.1006/icar.2000.6539。
\item Jean Meeus,天文算法(弗吉尼亚州里士满: willmann-bell，1998)。273页上九个极端的平均值。所有这些都在平均值的0.02 AU以内。
\item 凯撒，M。L.;Desch, M.D.;沃里克，J。W.;皮尔斯，J.B.(1980)。“旅行者探测到土星的非热无线电发射”。科学。209(4462):1238-40.Bibcode:1980Sci。209.1238K。doi:10.1126/science.209.4462.1238。高密度脂蛋白:2060/19800013712。PMID17811197。S2CID44313317。  
\item 本顿，朱利叶斯 (2006)。土星以及如何观察它。“天文学家观测指南” (第11版)。斯普林格科学与商业。第136页。ISBN 978-1-85233-887-9。已存档从原始的2017年2月16日。已检索2月9日2016。
\item “科学家发现土星的自转周期是一个谜”。NASA.2004年6月28日。已存档从2011年7月29日开始。已检索3月22日2007。
\item 弗雷泽凯恩 (2008年6月30日)。"土星"。今天的宇宙。已存档从2011年10月25日起。已检索8月17日2011年。
\item 安德森，J。D.;舒伯特，G.(2007)。“土星的引力场，内部旋转和内部结构”(PDF)。科学。317(5843):1384-1387。Bibcode:2007Sci。317.1384A。doi:10.1126/science.1144835。PMID17823351。S2CID19579769。存档自原来的(PDF)2020年4月12日。   
\item “土卫二间歇泉掩盖了土星日的长度”(新闻稿)。NASA喷气推进实验室2007年3月22日。存档自原来的2008年12月7日。已检索3月22日2007。
\item Gurnett, D.A.; 等人 (2007年)。“土星等离子盘内部区域的可变旋转周期”(PDF)。科学。316(5823):442-5.Bibcode:2007Sci...316 .. 442克。doi:10.1126/science.1138562。PMID17379775。S2CID46011210。存档自原来的(PDF)2020年2月12日。   
\item 巴格纳，F.(2007)。“土星自转的新自旋”。科学。316(5823):380-1.doi:10.1126/science.1142329。PMID17446379。S2CID118878929。  
\item Hui，Man-To; Wiegert，Paul A.; 等人 (2024年9月)。“2019 UO14: 土星的瞬态木马”。arXiv:2409.19725[astro-ph.EP]。
\item 侯，X。Y.; 等人 (2014年1月)。“土星木马: 动态的观点”。皇家天文学会月刊。437(2):1420-1433。Bibcode:2014MNRAS.437.1420H。doi:10.1093/mnras/stt1974。
\item 乔纳森·奥卡拉汉 (2025年3月11日)。“土星获得128颗新卫星，使其总数达到274颗”。纽约时报。已检索3月13日2025。
\item “Mpec 2025-e153: 61颗新的土星卫星”。
\item “Mpec 2025-e154: 三十四颗新的土星卫星”。
\item “Mpec 2025-e155: 三十三颗新的土星卫星”。
\item Tiscareno，Matthew (2013年7月17日)。“土星A环中螺旋桨的数量”。天文杂志。135(3):1083-1091。arXiv:0710.4547。Bibcode:2008AJ....135.1083T。doi:10.1088/0004-6256/135/3/1083。S2CID28620198。 
\item Brunier，Serge (2005)。太阳系航行。剑桥大学出版社。第164页。ISBN 978-0-521-80724-1。
\item 琼斯，G。等 (2008年3月7日)。“土星最大的冰冷卫星Rhea的尘埃晕”(PDF)。科学。319(5868):1380-1384。Bibcode:2008Sci...319.1380J。doi:10.1126/science.1151524。PMID18323452。S2CID206509814。存档自原来的(PDF)2018年3月8日   
\item 南希阿特金森 (2010年11月26日)。“在土星的月亮Rhea周围发现了微弱的氧气气氛”。今天的宇宙。已存档从2012年9月25日起。已检索7月20日2011年。
\item NASA (2010年11月30日)。稀薄的空气: 在土星的月亮Rhea上发现的氧气气氛。科学日报。已存档从2011年11月8日的原件。已检索7月23日2011年。
\item Ryan，Clare (2010年11月26日)。“卡西尼号揭示了土星月亮瑞亚的氧气气氛”。UCL Mullard空间科学实验室。已存档从2011年9月16日的原件。已检索7月23日2011年。
\item “已知的土星卫星”。地磁系。存档自原来的2011年9月26日。已检索6月22日2010。
\item “卡西尼号发现碳氢化合物降雨可能会填满泰坦湖”。科学日报。2009年1月30日已存档从2011年11月9日的原件。已检索7月19日2011年。
\item "旅行者-泰坦"。NASA喷气推进实验室2010年10月18日。存档自原来的2011年10月26日。已检索7月19日2011年。
\item “土卫六上碳氢化合物湖的证据”。NBC新闻。美联社。2006年7月25日。已存档从原件到2014年8月24日。已检索7月19日2011年。
\item “碳氢化合物湖终于在土卫六上确认”。《宇宙》杂志。2008年7月31日。存档自原来的2011年11月1日。已检索7月19日2011年。
\item 曼努埃尔·洛佩斯-普埃塔斯 (2013年6月6日)。“PAH在泰坦的高层大气中”。CSIC。已存档从2016年8月22日的原件。已检索6月6日2013。
\item 戴克斯，普雷斯顿; 等人 (2014年6月23日)。“泰坦的积木可能早于土星”。NASA。已存档从2018年9月9日的原件。已检索6月24日2014。
\item 巴特斯比，斯蒂芬 (2008年3月26日)。“土星的卫星土卫二令人惊讶地像彗星”。新科学家。已存档从原件到2015年6月30日。已检索4月16日2015。
\item NASA (2008年4月21日)。“土星的卫星土卫二上会有生命吗？”。科学日报。已存档从2011年11月9日的原件。已检索7月19日2011年。
\item 马德里加尔，亚历克西斯 (2009年6月24日)。“寻找土星月亮上的生命升温”。有线科学。已存档从2011年9月4日的原件。已检索7月19日2011年。
\item 斯波茨，彼得·N.(2005年9月28日)。“地球以外的生命？潜在的太阳系站点弹出“。今日美国。已存档从2008年7月26日开始。已检索7月21日2011年。
\item Pili，Unofre (2009年9月9日)。“土卫二: 土星的月亮，有水的液态海洋”。Scienceray。存档自原来的2011年10月7日。已检索7月21日2011年。
\item “最有力的证据表明土卫二隐藏着咸水海洋”。菲索格。2011年6月22日。已存档从2011年10月19日的原件。已检索7月19日2011年。
\item 考夫曼，马克 (2011年6月22日)。“土星的卫星土卫二显示出其表面下有海洋的证据”。华盛顿邮报。已存档从2012年11月12日的原件。已检索7月19日2011年。
\item Greicius，Tony; 等人 (2011年6月22日)。“卡西尼号在土星月亮上捕捉到海洋般的喷雾”。NASA.已存档从2011年9月14日的原件。已检索9月17日2011年。
\item 周，费利西亚; 戴奇，普雷斯顿; 韦弗，唐娜; 维拉德，雷 (2017年4月13日)。NASA任务为我们太阳系的 “海洋世界” 提供了新的见解。NASA.已存档从原件到2017年4月20日。已检索4月20日2017。
\item 普拉特，简; 等人 (2014年4月14日)。“NASA卡西尼号图像可能揭示土星卫星的诞生”。NASA。已存档从2019年4月10日的原件。已检索4月14日2014。
\item Poulet F.等人 (2002年)。“土星环的组成”。伊卡洛斯。160(2): 350。Bibcode:2002年Icar .. 160 .. 350便士。doi:10.1006/icar.2002.6967。已存档从2019年7月29日开始。已检索6月28日2019。
\item 卡罗琳·波科。“关于土星环的问题”。CICLOPS网站。已存档从2012年10月3日的原件。已检索6月18日2017。
\item Canup，罗宾·M。(2010年12月)。“通过从丢失的土卫六大小的卫星上进行质量清除，土星环和内部卫星的起源”。性质。468(7326):943-946。Bibcode:2010自然468.943C。doi:10.1038/自然09661。ISSN1476-4687。PMID21151108。已检索5月22日2024。  
\item 克里达，A.；Charnoz, S.(2012年11月30日)。“从太阳系古代大质量环形成常规卫星”。科学。338(6111):1196-1199。arXiv:1301.3808。Bibcode:2012Sci。338.1196C。doi:10.1126/science.1226477。PMID23197530。 
\item 沙诺兹，塞巴斯蒂安; 莫比德利，亚历山德罗; 唐斯，卢克; 鲑鱼，朱利安 (2009年2月)。“土星的光环是在后期的猛烈轰炸中形成的吗？”。伊卡洛斯。199(2):413-428。arXiv:0809.5073。Bibcode:2009年Icar .. 199 .. 413C。doi:10.1016/j.icarus.2008.10.019。已检索5月22日2024。
\item 肯普夫，萨沙; 阿尔托比利，尼古拉斯; 施密特，于尔根; 库齐，杰弗里·N；埃斯特拉达，保罗·R；拉尔夫·斯拉马 (2023年5月12日)。“微流星体进入土星环将它们的年龄限制在不超过几亿年”。科学进展。9(19) eadf8537。Bibcode:2023年SciA .... 9F8537K。doi:10.1126/sciadv.adf8537。PMC10181170。PMID37172091。  
\item 安德鲁，罗宾·乔治 (2023年9月28日)。“土星环可能是在最近一次令人惊讶的两颗卫星撞击中形成的 -- 研究人员完成了一个复杂的模拟，支持这颗巨大行星的珠宝出现在数亿年前的想法，不是数十亿”。纽约时报。已存档从2023年9月29日开始。已检索9月29日2023。
\item Teodoro，L.F.A.; 等人 (2023年9月27日)。“土星环和中型卫星的最新影响起源”。天体物理学杂志。955(2): 137。arXiv:2309.15156。Bibcode:2023年4月。955 .. 137吨。doi:10.3847/1538-4357/acf4ed。
\item 智慧，杰克; Dbouk，Rola; Militzer，Burkhard; 哈伯德，威廉B.；弗朗西斯·尼莫; 布兰娜·G·唐尼；法语，理查德·G。(2022年9月16日)。“失去一颗卫星可以解释土星的倾斜和年轻的光环”。科学。377(6612):1285-1289。Bibcode:2022Sci...377.1285W。doi:10.1126/科学.abn1234。高密度脂蛋白:1721.1/148216。PMID36107998。S2CID252310492。  
\item Cowen，Rob (1999年11月7日)。“发现最大的已知行星环”。科学新闻。已存档从2011年8月22日的原件。已检索4月9日2010。
\item Russell，Randy (2004年6月7日)。“土星的卫星和环”。宇宙的窗户已存档从2011年9月4日的原件。已检索7月19日2011年。
\item NASA喷气推进实验室 (2005年3月3日)。“NASA的卡西尼号飞船继续取得新发现”。科学日报。已存档从2011年11月8日的原件。已检索7月19日2011年。
\item “观察土星”。国家海事博物馆。2015年8月20日存档自原来的2007年4月22日。已检索7月6日2007。
\item 萨克斯，A。(1974年5月2日)。“巴比伦观测天文学”。伦敦皇家学会的哲学交易。276(1257):43-50。Bibcode:1974RSPTA.276...43S。doi:10.1098/rsta.1974.0008。JSTOR74273。S2CID121539390。  
\item Φαίνων。亨利·乔治·利德尔;罗伯特·斯科特;中级希腊英语词典在Perseus项目。
\item 西塞罗，De Natura Deorum。
\item Neuh ä user，R; Torres，G; Mugrauer，M; Neuh ä user，D L; Chapman，J; Luge，D; Cosci，M (2022年10月11日)。“两千年来，从历史记录中得出的Betelgeuse和Antares的颜色演变，是对质量和年龄的新限制”。皇家天文学会月刊。516(1):693-719。arXiv:2207.04702。doi:10.1093/mnras/stac1969。ISSN0035-8711。 
\item 《星夜时报》。Imaginova公司，2006年。存档自原来的2009年10月1日。已检索7月5日2007。
 “行星的希腊名字”。2010年4月25日。已存档从2010年5月9日的原件。已检索7月14日2012。土星的希腊名字是Kronos。泰坦克洛诺斯是宙斯,而土星是罗马的农业之神。另请参见希腊关于地球的文章。
 邦尼尔公司 (1893年4月)。“流行杂项-关于土星的迷信”。《科普月刊》: 862。已存档从2017年2月17日的原件。已检索2月9日2016。
 Jan Jakob Maria De Groot (1912)。中国的宗教: 大学主义。道家与儒学研究的关键。美国关于宗教史的讲座。第10卷。G.P。普特南的儿子。第300页。已存档从2011年7月22日开始。已检索1月8日2010。
 克鲁普，托马斯 (1992)。日本的数字游戏: 现代日本对数字的使用和理解。Routledge. pp. 39-40.ISBN 978-0-415-05609-0。
 Hulbert，Homer Bezaleel (1909)。韩国的通过。双日，佩奇和公司。p。 426。已检索1月8日2010。
 艾比塞斯纳 (2009年11月15日)。“土星是什么时候被发现的？”。今天的宇宙。已存档从2012年2月14日的原件。已检索7月21日2011年。
 “法师，第一本书: 天体情报员: 第二十八章”。Sacred-Text.com。已存档从原件到2018年6月19日。已检索8月4日2018。
 “神话中的土星”。CrystalLinks.com。已存档从原来的2018年7月8日。已检索8月5日2018。
 凯瑟琳·拜尔 (2017年3月8日)。"行星精神印记-01土星精神"。ThoughtCo.com。已存档从2018年8月4日的原件。已检索8月3日2018。
 "含义和起源: Zazel"。FamilyEducation.com。2014。已存档从原来的2015年1月2日。已检索8月3日2018。拉丁语: 天使因爱情召唤而召唤
 “天使的存在”。Hafapea.com。1998。存档自原来的2018年7月22日。已检索8月3日2018。所罗门的爱之天使仪式
 伊士曼，杰克 (1998)。“双筒望远镜中的土星”。丹佛天文学会。存档自原来的2011年7月28日。已检索9月3日2008。
 陈，加里 (2000)。“土星: 历史时间表”。已存档从2011年7月16日开始。已检索7月16日2007。
 该隐，弗雷泽 (2008年7月3日)。“土星的历史”。今天的宇宙。已存档从2012年1月26日开始。已检索7月24日2011年。
 该隐，弗雷泽 (2008年7月7日)。“关于土星的有趣事实”。今天的宇宙。已存档从2011年9月25日起。已检索9月17日2011年。
 该隐，弗雷泽 (2009年11月27日)。“谁发现了土星？”。今天的宇宙。已存档从原件到2012年7月18日。已检索9月17日2011年。
 米克，凯瑟琳。“土星: 发现的历史”。存档自原来的2011年7月23日。已检索7月15日2007。
 塞缪尔·G·巴顿.(1946年4月)。“卫星的名字”。热门天文学。第54卷。pp。 122-130。Bibcode:1946帕…… 54..122b。
 柯伊伯，杰拉德P.(1944年11月)。“泰坦: 一颗有大气层的卫星”。天体物理学杂志。100:378-388。Bibcode:1944年4月。100 .. 378K。doi:10.1086/144679。
 “先锋10号和11号飞船”。任务描述。存档自原来的2006年1月30日。已检索7月5日2007。
 “土星任务”。行星社会。2007。已存档从2011年7月28日起。已检索7月24日2007。
 Spence，Pam (1999)。揭示了宇宙。剑桥大学出版社。第64页。ISBN 978-0-521-64239-2。
 戴奇，普雷斯顿; 等人 (2014年7月28日)。“卡西尼号飞船在冰冷的土星月球上揭示了101个间歇泉和更多”。NASA。已存档从原始的2017年7月14日。已检索7月29日2014。
 Lebreton，jean-pierre; 等人 (2005年12月)。“惠更斯探测器在土卫六上的下降和着陆概述”。性质。438(7069):758-764。Bibcode:2005 Natur.438..758L。doi:10.1038/nature04347。PMID16319826。S2CID4355742。  
 “天文学家在土星发现了巨大的闪电风暴”。科学日报有限责任公司。2007。已存档从2011年8月28日的原件。已检索7月27日2007。
 彭斯，迈克尔 (2006年3月9日)。“NASA的卡西尼号在土卫二上发现了潜在的液态水”。NASA喷气推进实验室。已存档从2011年8月11日的原件。已检索6月3日2011年。
 理查德·洛维特.(2011年5月31日)。“土卫二被命名为外星生命最甜蜜的地方”。性质news.2011.337.doi:10.1038/news.2011.337。已存档从2011年9月5日的原件。已检索6月3日2011年。
 凯西喀山 (2011年6月2日)。“土星的土卫二移至“ 最有可能拥有生命 ”列表的顶部”。每日银河已存档从2011年8月6日的原件。已检索6月3日2011年。
 大卫，志贺 (2007年9月20日)。“在土星周围发现了微弱的新环”。NewScientist.com。已存档从2008年5月3日的原件。已检索7月8日2007。
 Rincon，Paul (2007年3月14日)。“探测器揭示了土星月球上的海洋”。英国广播公司。已存档从2011年11月11日的原件。已检索9月26日2007。
 Rincon，Paul (2006年11月10日)。“巨大的 '飓风' 肆虐土星”。英国广播公司。已存档从2011年9月2日的原件。已检索7月12日2007。
 “任务概述-简介”。卡西尼冬至任务。NASA / JPL。2010。存档自原来的2011年8月7日。已检索11月23日2010。
 “土星北极的大风暴”。3 News NZ。2013年4月30日已存档从原件到2014年7月19日。已检索4月30日2013。
 布朗，德韦恩; 坎蒂洛，劳里; 戴奇，普雷斯顿 (2017年9月15日)。美国宇航局的卡西尼号飞船结束了对土星的历史性探索。NASA。已存档从2019年5月9日的原始。已检索9月15日2017。
 Chang，Kenneth (2017年9月14日)。卡西尼号消失在土星，它的使命庆祝和哀悼。纽约时报。已存档从原来的2018年7月8日。已检索9月15日2017。
 杰夫·福斯特 (2016年1月8日)。“NASA扩大了下一个新前沿竞赛的前沿”。SpaceNews。已存档从原始的2017年8月18日。已检索4月20日2017。
 Nola Taylor Redd (2017年4月25日)。"“蜻蜓” 无人机可以探索土星卫星泰坦。Space.com。已存档从原件到2019年6月30日。已检索6月13日2020。
 "土星环边缘"。古典天文学。2013。存档自原来的2013年11月5日。已检索8月4日2013。
 理查德·W·施穆德.小 (2003年冬季)。“2002-03年的土星”。佐治亚科学杂志。61(4)。ISSN0147-9369。已存档从原来的2015年9月24日。已检索6月29日2015。 
 Tanya Hill; 等人 (2014年5月9日)。“无论如何，明亮的土星将在澳大利亚闪烁一个小时”。对话。已存档从原来的2014年5月10日。已检索5月11日2014。
 麦金尼，理查德L.(2005)。“木星和外行星”。在加里·韦斯特法尔(ed.)。格林伍德科幻小说和幻想百科全书: 主题，作品和奇迹。格林伍德出版集团。第449页。ISBN 978-0-313-32951-7。最早以土星为主题的小说可能是伏尔泰的Micromégas(1750)。很久以后，土星成为Poul Anderson的 “土星游戏” (1981年) 和Michael A的中心。麦科勒姆的土星的云(1991)，人类城市漂浮在土星的大气层中。地球的大气层也是罗伯特·F·两脑四公里宽的生物的家园。前锋的土星Rukh(1997)。土星最大的卫星，泰坦-有趣的是它的厚厚的大气层-在艾伦·诺斯1954年的少年小说中被殖民，泰坦上的麻烦,而斯蒂芬·巴克斯特的泰坦(1997) 是关于卫星的太空任务。
 加里·韦斯特法尔(2021)。"土星"。历史上的科幻文学: 百科全书。ABC-CLIO. pp. 553-555。ISBN 978-1-4408-6617-3。
 布莱恩·斯坦福德(2006)。"土星"。科学事实与科幻小说: 百科全书。泰勒和弗朗西斯。pp. 458-459。ISBN 978-0-415-97460-8。
\end{enumerate}










